% !TEX program = xelatex
\documentclass[11pt,a4paper]{article}
\usepackage[margin=2.2cm,top=2cm]{geometry}
\usepackage{fontspec}
\usepackage{polyglossia}
\setdefaultlanguage{turkish}
\usepackage{amsmath,amssymb}
\usepackage{graphicx}
\usepackage{booktabs}
\usepackage{enumitem}
\usepackage{multicol}
\usepackage[most]{tcolorbox}
\usepackage{fancyhdr}
\usepackage{titlesec}
\usepackage{xcolor}
\usepackage{needspace}
\usepackage[hidelinks]{hyperref}

\definecolor{anamavi}{RGB}{0,45,98}
\definecolor{cozumyesil}{RGB}{0,110,50}
\definecolor{uyarikirmizi}{RGB}{180,20,20}

\titleformat{\section}{\large\bfseries\color{anamavi}}{\thesection.}{0.5em}{}
\titleformat{\subsection}{\normalsize\bfseries\color{anamavi}}{\thesubsection}{0.5em}{}

\pagestyle{fancy}\fancyhf{}
\lhead{\small MMAT105 Genel Matematik}
\rhead{\small Hafta 12 -- Ders Notu}
\cfoot{\small\thepage}
\renewcommand{\headrulewidth}{0.4pt}

\newtcolorbox{kural}[1][Kural]{colback=anamavi!5,colframe=anamavi,title=#1,fonttitle=\bfseries,breakable,left=4pt,right=4pt,top=2pt,bottom=2pt}
\newtcolorbox{dikkat}[1][Dikkat]{colback=uyarikirmizi!5,colframe=uyarikirmizi,title=#1,fonttitle=\bfseries,breakable,left=4pt,right=4pt,top=2pt,bottom=2pt}
\newcounter{ornek}
\newenvironment{ornek}{\refstepcounter{ornek}\par\medskip\needspace{7\baselineskip}\noindent\textbf{\color{anamavi}Örnek \theornek.}\ }{\par}
\newenvironment{cozum}{\begin{tcolorbox}[colback=cozumyesil!6,colframe=cozumyesil,boxrule=0.5pt,left=4pt,right=4pt,top=2pt,bottom=2pt,breakable]\textbf{\color{cozumyesil}Çözüm.}\ }{\end{tcolorbox}}
\newcounter{alis}
\newenvironment{alistirma}{\refstepcounter{alis}\par\medskip\needspace{9\baselineskip}\noindent\textbf{\color{anamavi}Alıştırma \thealis.}\ }{\par}
\newcommand{\kendinizoncesi}{\needspace{12\baselineskip}}
\newcommand{\ab}[1]{\overline{#1}}
\setlength{\parindent}{0pt}
\setlength{\parskip}{3pt}

\begin{document}

% ---------------- Başlık ----------------
\begin{minipage}[c]{\textwidth}
{\LARGE\bfseries\color{anamavi} MMAT105 Genel Matematik}\\[4pt]
{\large Hafta 12 Ders Notu: Oran ve Orantı (Doğru, Ters,\\ Bileşik Orantı) ve Ortalamalar}\\[6pt]
{\small 2026--2027 Güz}
\end{minipage}

\vspace{4pt}\noindent\textcolor{anamavi}{\rule{\textwidth}{1pt}}

\textbf{Bu notun kullanımı.} Her bölümde önce kısa bir kural kutusu, ardından çözümlü örnekler yer alır. Örnekleri çözümüne bakmadan önce kendiniz deneyin. Notun sonunda önce çözümlü alıştırmalar, ardından cevap anahtarlı ``Kendiniz Çözün'' alıştırmaları yer alır.

% =====================================================
\section{Oran ve Orantı}

\begin{kural}[Tanımlar]
\textbf{Oran:} İki çokluğun bölümüdür: \(\dfrac{a}{b}\) \ (\(b\neq 0\)). Aynı birimli çoklukların oranı birimsizdir; farklı birimli çoklukların oranı birimlidir (ör. hız, km/sa).\\
\textbf{Orantı:} İki oranın eşitliğidir. \(b,d\neq 0\) için
\[\frac{a}{b}=\frac{c}{d}\ \Longleftrightarrow\ a\cdot d=b\cdot c\qquad(\text{içler--dışlar çarpımı})\]
\end{kural}

\begin{dikkat}
Farklı birimli çokluklar önce aynı birime çevrilir (\(1\) saat \(=60\) dakika, \(1\) m \(=100\) cm). \(\dfrac{a}{b}=\dfrac{2}{3}\) ise \(a=2\), \(b=3\) olmak zorunda değildir; \(a=2k\), \(b=3k\) (\(k\neq 0\)) yazılır.
\end{dikkat}

\begin{ornek}
\(\dfrac{3x-2}{x+1}=\dfrac{4}{3}\) olduğuna göre \(x\) kaçtır?
\end{ornek}
\begin{cozum}
\(3(3x-2)=4(x+1)\ \Rightarrow\ 9x-6=4x+4\ \Rightarrow\ 5x=10\ \Rightarrow\ x=2\). \quad Kontrol: \(\dfrac{4}{3}=\dfrac{4}{3}\) \checkmark
\end{cozum}

\begin{ornek}
\(\dfrac{a}{b}=\dfrac{5}{2}\) olduğuna göre \(\dfrac{a+2b}{a-b}\) kaçtır?
\end{ornek}
\begin{cozum}
\(a=5k,\ b=2k\): \quad \(\dfrac{5k+4k}{5k-2k}=\dfrac{9k}{3k}=3\).
\end{cozum}

\begin{ornek}
\(\dfrac{2a-b}{a+b}=\dfrac{1}{4}\) olduğuna göre \(\dfrac{a}{b}\) kaçtır?
\end{ornek}
\begin{cozum}
\(4(2a-b)=a+b\ \Rightarrow\ 8a-4b=a+b\ \Rightarrow\ 7a=5b\ \Rightarrow\ \dfrac{a}{b}=\dfrac{5}{7}\).\\
Kontrol: \(a=5,\ b=7\): \(\dfrac{10-7}{12}=\dfrac{3}{12}=\dfrac{1}{4}\) \checkmark
\end{cozum}

% =====================================================
\section{Orantı Sabiti ve Paylaştırma}

\begin{kural}[Orantı sabiti]
\(\dfrac{a}{b}=\dfrac{c}{d}=\dfrac{e}{f}=k\) (paydalar sıfırdan farklı) ise \(a=bk,\ c=dk,\ e=fk\) ve
\[\frac{a+c+e}{b+d+f}=k,\qquad \frac{m\,a+n\,c}{m\,b+n\,d}=k\quad(mb+nd\neq 0).\]
\(\dfrac{a}{2}=\dfrac{b}{3}=\dfrac{c}{5}=k\) ise \(a=2k,\ b=3k,\ c=5k\); bu, \(a:b:c=2:3:5\) demektir.\\
\textbf{Paylaştırma:} Bir miktar \(2,3,5\) ile (doğru) orantılı paylaştırılırsa paylar \(2k,3k,5k\) olur ve \(2k+3k+5k=\) miktar. \(2,3,4\) ile \textbf{ters orantılı} paylaştırılırsa paylar \(\frac12,\frac13,\frac14\) ile doğru orantılıdır; paydaların EKOK'u ile çarpılarak \(6:4:3\) elde edilir.
\end{kural}

\begin{ornek}
\(\dfrac{a}{3}=\dfrac{b}{5}=\dfrac{c}{7}\) ve \(2a+b-c=16\) olduğuna göre \(a+b+c\) kaçtır?
\end{ornek}
\begin{cozum}
\(a=3k,\ b=5k,\ c=7k\): \ \(6k+5k-7k=4k=16\Rightarrow k=4\). \quad \(a+b+c=15k=60\) \ (\(a=12,\ b=20,\ c=28\)).
\end{cozum}

\begin{ornek}
\(a:b=3:4\) ve \(b:c=6:5\) olduğuna göre \(a:b:c\) kaçtır? \(a+b+c=62\) ise \(c\) kaçtır?
\end{ornek}
\begin{cozum}
Ortak terim \(b\): \(\mathrm{EKOK}(4,6)=12\). \ \(a:b=9:12\), \ \(b:c=12:10\) \(\Rightarrow\) \(a:b:c=9:12:10\).\\
\(31k=62\Rightarrow k=2\) \(\Rightarrow\) \(c=20\) \ (\(a=18,\ b=24\)).
\end{cozum}

\begin{ornek}
\(360\) TL, \(3\), \(4\) ve \(5\) ile doğru orantılı paylaştırılıyor. Ayrıca \(390\) TL, \(2\), \(3\) ve \(4\) ile ters orantılı paylaştırılıyor. Her iki paylaştırmada payları bulunuz.
\end{ornek}
\begin{cozum}
Doğru orantı: \(3k+4k+5k=12k=360\Rightarrow k=30\) \(\Rightarrow\) \(90,\ 120,\ 150\) TL.\\
Ters orantı: \(\tfrac12:\tfrac13:\tfrac14=6:4:3\) (EKOK \(12\) ile çarpıldı). \(13k=390\Rightarrow k=30\) \(\Rightarrow\) \(180,\ 120,\ 90\) TL.\\
Dikkat: ters orantıda en büyük sayıya (\(4\)) en küçük pay (\(90\) TL) düşer.
\end{cozum}

% =====================================================
\section{Harita Ölçeği}

\begin{kural}[Ölçek]
\[\text{Ölçek}=\frac{\text{Haritadaki uzunluk}}{\text{Gerçek uzunluk}}\quad(\text{aynı birimle}).\qquad 1\text{ km}=1000\text{ m}=100\,000\text{ cm}.\]
\(1:2\,500\,000\) ölçeğinde haritadaki \(1\) cm, gerçekte \(2\,500\,000\) cm \(=25\) km'dir.
\end{kural}

\begin{ornek}
\(1:2\,500\,000\) ölçekli bir haritada iki şehir arası \(4{,}8\) cm'dir. Gerçek uzaklık kaç km'dir? \(1:200\) ölçekli bir planda bir odanın kenarları \(5\) cm ve \(3{,}5\) cm'dir. Odanın gerçek kenar uzunlukları kaç m'dir?
\end{ornek}
\begin{cozum}
\(4{,}8\cdot 2\,500\,000=12\,000\,000\) cm \(=120\) km.\\
\(5\cdot 200=1000\) cm \(=10\) m; \quad \(3{,}5\cdot 200=700\) cm \(=7\) m.
\end{cozum}

% =====================================================
\section{Doğru ve Ters Orantı}

\begin{kural}[Doğru ve ters orantı]
\(x\) ve \(y\) pozitif çokluklar, \(k>0\) bir sabit olsun.\\
\textbf{Doğru orantı:} \(\dfrac{y}{x}=k\), yani \(y=kx\). \(x\) kaç katına çıkarsa \(y\) de o kata çıkar. \ \(\dfrac{y_1}{x_1}=\dfrac{y_2}{x_2}\).\\
\textbf{Ters orantı:} \(x\cdot y=k\). \(x\) kaç katına çıkarsa \(y\) o kadar kata bölünür. \ \(x_1y_1=x_2y_2\).\\
Tipik ters orantılar (aynı iş / aynı yol / aynı havuz): işçi sayısı -- gün, hız -- zaman, musluk sayısı -- süre, dişli çarkta diş sayısı -- tur sayısı.
\end{kural}

\begin{dikkat}
``Biri artarken diğeri de artıyor'' doğru orantı için yeterli değildir; \textbf{oranın sabit} olması gerekir. Örneğin \(y=2x+1\) artandır ama \(\frac{y}{x}\) sabit olmadığından doğru orantı değildir.
\end{dikkat}

\begin{ornek}
\(y\), \(x\) ile ters orantılıdır. \(x=3\) iken \(y=20\)'dir. \(y=12\) iken \(x\) kaçtır?
\end{ornek}
\begin{cozum}
\(x\cdot y=3\cdot 20=60\) sabit. \quad \(x\cdot 12=60\Rightarrow x=5\).
\end{cozum}

\begin{ornek}
Bir çiftlikteki \(40\) ineğe \(30\) gün yetecek kadar yem vardır. \(10\) gün sonra \(20\) inek satılıyor. Kalan yem, kalan ineklere kaç gün yeter?
\end{ornek}
\begin{cozum}
\(10\) gün sonra kalan yem \(40\) ineğe \(30-10=20\) gün yeterdi. Kalan inek: \(40-20=20\).\\
İnek sayısı ile gün ters orantılı: \(40\cdot 20=20\cdot x\Rightarrow x=40\) gün.
\end{cozum}

\begin{ornek}
Bir izci grubuna \(20\) gün yetecek kadar yiyecek vardır. \(5\) gün sonra gruba \(10\) izci daha katılıyor ve kalan yiyecek tüm gruba \(10\) gün yetiyor. Başlangıçta grupta kaç izci vardı?
\end{ornek}
\begin{cozum}
Başlangıçta \(n\) izci. \(5\) gün sonra kalan yiyecek \(n\) izciye \(15\) gün yeterdi. Kişi sayısı ile gün ters orantılı:
\[n\cdot 15=(n+10)\cdot 10\ \Rightarrow\ 15n=10n+100\ \Rightarrow\ n=20\]
\end{cozum}

% =====================================================
\section{Bileşik Orantı}

\begin{kural}[Bileşik orantı]
\(a\); \(b\) ile doğru, \(c\) ile ters orantılı ise \(a=k\cdot\dfrac{b}{c}\), yani \(\dfrac{a\cdot c}{b}=k\) sabittir.\\
Aynı verimle çalışanlar için:
\[\frac{\text{İş miktarı}}{\text{İşçi sayısı}\cdot\text{Günlük saat}\cdot\text{Gün sayısı}}=\text{sabit}\]
\end{kural}

\begin{ornek}
\(x\) sayısı \(y\) ile ters, \(z\) ile doğru orantılıdır. \(y=4\), \(z=6\) iken \(x=9\)'dur. \(y=6\), \(z=10\) iken \(x\) kaçtır?
\end{ornek}
\begin{cozum}
\(x=k\cdot\dfrac{z}{y}\Rightarrow k=\dfrac{x\cdot y}{z}=\dfrac{9\cdot 4}{6}=6\). \quad \(x=6\cdot\dfrac{10}{6}=10\).
\end{cozum}

\begin{ornek}
\(12\) işçi günde \(8\) saat çalışarak bir işi \(15\) günde bitiriyor. Aynı nitelikteki \(9\) işçi günde \(10\) saat çalışarak aynı işi kaç günde bitirir?
\end{ornek}
\begin{cozum}
İş aynı olduğundan işçi \(\cdot\) saat \(\cdot\) gün sabittir:
\[12\cdot 8\cdot 15=9\cdot 10\cdot x\ \Rightarrow\ 1440=90x\ \Rightarrow\ x=16\text{ gün}\]
\end{cozum}

% =====================================================
\section{Ortalamalar}

\begin{kural}[Aritmetik ve ağırlıklı ortalama]
\(\text{AO}=\dfrac{x_1+x_2+\cdots+x_n}{n}\), \quad Toplam \(=n\cdot\text{AO}\).\\
Bir sayı eklenir veya çıkarılırsa: yeni ortalama \(=\dfrac{\text{eski toplam}\pm\text{sayı}}{\text{yeni sayı adedi}}\).\\
\textbf{Ağırlıklı ortalama:} \(\dfrac{n_1\,o_1+n_2\,o_2}{n_1+n_2}\) (\(n_i\): grup büyüklüğü ya da notun ağırlığı).
\end{kural}

\begin{kural}[{Geometrik ve harmonik ortalama; ardışık sayılar}]
İki pozitif sayı \(a,b\) için: \ \(\text{GO}=\sqrt{ab}\), \ \(\text{HO}=\dfrac{2ab}{a+b}\), \ \(\text{AO}\ge\text{GO}\ge\text{HO}\) (eşitlik yalnızca \(a=b\) iken).\\
Ortalama hız \(=\dfrac{\text{toplam yol}}{\text{toplam zaman}}\); aynı yol \(v_1\) ile gidilip \(v_2\) ile dönülürse \(\dfrac{2v_1v_2}{v_1+v_2}\).\\
Ardışık (eşit aralıklı) sayıların ortalaması \(=\dfrac{\text{ilk}+\text{son}}{2}\); terim sayısı tek ise ortanca terim.
\end{kural}

\begin{dikkat}
Farklı büyüklükteki grupların ortak ortalaması, ortalamaların ortalaması değildir. Gidiş--dönüşte ortalama hız \(\frac{v_1+v_2}{2}\) değildir.
\end{dikkat}

\begin{ornek}
Aritmetik ortalaması \(20\) olan \(6\) sayıdan biri çıkarılınca kalan sayıların ortalaması \(18\) oluyor. Çıkarılan sayı kaçtır?
\end{ornek}
\begin{cozum}
Eski toplam: \(6\cdot 20=120\); \ yeni toplam: \(5\cdot 18=90\). \quad Çıkarılan sayı: \(120-90=30\).
\end{cozum}

\begin{ornek}
Bir derste vize notunun ağırlığı \(1\), final notunun ağırlığı \(2\)'dir. Vizeden \(48\) alan bir öğrencinin ağırlıklı ortalamasının \(60\) olması için finalden kaç alması gerekir?
\end{ornek}
\begin{cozum}
\(\dfrac{1\cdot 48+2\cdot f}{1+2}=60\Rightarrow 48+2f=180\Rightarrow 2f=132\Rightarrow f=66\).
\end{cozum}

\begin{ornek}
\(30\) kişilik bir sınıfın not ortalaması \(62\)'dir. Sınıftaki \(12\) kızın ortalaması \(68\) ise erkeklerin ortalaması kaçtır?
\end{ornek}
\begin{cozum}
Sınıf toplamı: \(30\cdot 62=1860\); \ kızların toplamı: \(12\cdot 68=816\).\\
Erkek sayısı \(18\), toplamları \(1860-816=1044\). \quad Ortalama: \(\dfrac{1044}{18}=58\).
\end{cozum}

\begin{ornek}
\(2\) ve \(8\) sayılarının aritmetik, geometrik ve harmonik ortalamalarını bulunuz. \(\text{AO}\cdot\text{HO}\) çarpımını \(\text{GO}^2\) ile karşılaştırınız.
\end{ornek}
\begin{cozum}
\(\text{AO}=5\), \ \(\text{GO}=\sqrt{16}=4\), \ \(\text{HO}=\dfrac{2\cdot 16}{10}=3{,}2\). \ \(5>4>3{,}2\).\\
\(\text{AO}\cdot\text{HO}=5\cdot 3{,}2=16=\text{GO}^2\). (İki pozitif sayı için her zaman \(\frac{a+b}{2}\cdot\frac{2ab}{a+b}=ab\).)
\end{cozum}

\begin{ornek}
Bir araç A'dan B'ye saatte \(45\) km hızla gidip aynı yoldan saatte \(90\) km hızla dönüyor. Ortalama hızı kaç km/sa'tir?
\end{ornek}
\begin{cozum}
\(\dfrac{2\cdot 45\cdot 90}{45+90}=\dfrac{8100}{135}=60\) km/sa.\\
Kontrol: yol \(90\) km olsun: gidiş \(2\) sa, dönüş \(1\) sa; \ \(\dfrac{180}{3}=60\). (\(67{,}5\) değil.)
\end{cozum}

\begin{ornek}
Ardışık \(9\) tam sayının ortalaması \(23\)'tür. En küçüğü ve en büyüğü kaçtır? \(20\)'den \(60\)'a kadar (dahil) çift sayıların toplamı kaçtır?
\end{ornek}
\begin{cozum}
Ortanca (5. terim) \(23\) \(\Rightarrow\) en küçük \(23-4=19\), en büyük \(23+4=27\).\\
Çift sayılar: terim sayısı \(\dfrac{60-20}{2}+1=21\), ortalama \(\dfrac{20+60}{2}=40\), toplam \(21\cdot 40=840\).
\end{cozum}

% =====================================================
\section{Özet: Formül Kartı}
\begin{center}
\renewcommand{\arraystretch}{1.6}
\begin{tabular}{@{}ll@{}}
\toprule
Konu & Formül / Kural\\\midrule
Oran & \(\frac{a}{b}\), aynı birimle; birimsiz\\
Orantı & \(\frac{a}{b}=\frac{c}{d}\Leftrightarrow ad=bc\)\\
Orantı sabiti & \(\frac{a}{b}=\frac{c}{d}=\frac{e}{f}=k\Rightarrow \frac{a+c+e}{b+d+f}=k,\ \ \frac{ma+nc}{mb+nd}=k\)\\
Paylaştırma & \(2,3,5\) ile orantılı: \(2k,3k,5k\); \ ters orantılı: \(\frac12,\frac13,\frac15\) ile orantılı\\
Ölçek & haritadaki uzunluk / gerçek uzunluk; \ \(1\) km \(=100\,000\) cm\\
Doğru orantı & \(\frac{y}{x}=k\), \ \(\frac{y_1}{x_1}=\frac{y_2}{x_2}\)\\
Ters orantı & \(xy=k\), \ \(x_1y_1=x_2y_2\) \ (işçi--gün, hız--zaman, musluk--süre)\\
Bileşik orantı & \(a\); \(b\) ile doğru, \(c\) ile ters \(\Rightarrow a=k\frac{b}{c}\); \ \(\frac{\text{İş}}{n\cdot s\cdot g}\) sabit\\
Aritmetik ort. & \(\text{AO}=\frac{\text{toplam}}{n}\), \ toplam \(=n\cdot\text{AO}\)\\
Ağırlıklı ort. & \(\frac{n_1o_1+n_2o_2}{n_1+n_2}\)\\
GO, HO & \(\sqrt{ab}\), \ \(\frac{2ab}{a+b}\); \ \(\text{AO}\ge\text{GO}\ge\text{HO}\)\\
Ortalama hız & toplam yol / toplam zaman; gidiş--dönüş: \(\frac{2v_1v_2}{v_1+v_2}\)\\
Ardışık sayılar & ortalama \(=\frac{\text{ilk}+\text{son}}{2}\); terim sayısı \(\frac{\text{son}-\text{ilk}}{\text{artış}}+1\)\\\bottomrule
\end{tabular}
\end{center}

% =====================================================
\newpage
% Hafta 12 -- alıştırma gövdesi (ders notu ve ayrı alıştırma kağıdı tarafından \input edilir)
\section{Çözümlü Alıştırmalar}
Her soruyu önce kendiniz deneyin, sonra çözümü okuyun.

\begin{alistirma}
\(\dfrac{a}{b}=\dfrac{4}{7}\) ve \(a+b=44\) olduğuna göre \(b-a\) kaçtır?
\end{alistirma}
\begin{cozum}
\(a=4k\), \(b=7k\) yazalım: \ \(4k+7k=11k=44\Rightarrow k=4\).\\
\(a=16\), \(b=28\) \(\Rightarrow\) \(b-a=12\) \ (kısa yol: \(b-a=3k=12\)).
\end{cozum}

\begin{alistirma}
\(\dfrac{x}{4}=\dfrac{y}{5}=\dfrac{z}{6}\) ve \(x+y+z=45\) olduğuna göre \(x\cdot z\) kaçtır?
\end{alistirma}
\begin{cozum}
\(x=4k,\ y=5k,\ z=6k\): \ \(15k=45\Rightarrow k=3\).\\
\(x=12\), \(z=18\) \(\Rightarrow\) \(x\cdot z=216\).
\end{cozum}

\begin{alistirma}
\(2a=5b\) ve \(a-b=18\) olduğuna göre \(a\) ve \(b\) kaçtır?
\end{alistirma}
\begin{cozum}
\(2a=5b=10k\) diyelim (\(\mathrm{EKOK}(2,5)=10\)): \ \(a=5k\), \(b=2k\).\\
\(a-b=3k=18\Rightarrow k=6\) \(\Rightarrow\) \(a=30\), \(b=12\). \ Kontrol: \(2\cdot 30=60=5\cdot 12\) \checkmark
\end{cozum}

\begin{alistirma}
\(2600\) TL, üç kardeş arasında \(\dfrac{1}{2}\), \(\dfrac{1}{3}\) ve \(\dfrac{1}{4}\) sayıları ile orantılı olarak paylaştırılıyor. Her kardeş kaç TL alır?
\end{alistirma}
\begin{cozum}
Paydaların EKOK'u \(12\) ile çarpalım: \(\tfrac12:\tfrac13:\tfrac14=6:4:3\).\\
\(6k+4k+3k=13k=2600\Rightarrow k=200\) \(\Rightarrow\) paylar \(1200\), \(800\), \(600\) TL.\\
(Bu, \(2,3,4\) ile ters orantılı paylaştırma ile aynıdır.)
\end{cozum}

\begin{alistirma}
\(9\) işçi bir işi \(20\) günde bitirebiliyor. \(5\) gün çalıştıktan sonra aynı nitelikte \(6\) işçi daha katılıyor. Kalan iş kaç günde biter?
\end{alistirma}
\begin{cozum}
\(9\) işçi \(5\) gün çalıştı; kalan iş \(9\) işçinin \(20-5=15\) günlük işidir.\\
Kalan iş için işçi sayısı ile gün ters orantılı: \(9\cdot 15=(9+6)\cdot x\Rightarrow 135=15x\Rightarrow x=9\) gün.
\end{cozum}

\begin{alistirma}
Özdeş \(3\) makine günde \(8\) saat çalışarak \(6\) günde \(1800\) parça üretiyor. Bu makinelerden \(4\) tanesi günde \(6\) saat çalışarak \(2400\) parçayı kaç günde üretir?
\end{alistirma}
\begin{cozum}
\[\frac{1800}{3\cdot 8\cdot 6}=\frac{2400}{4\cdot 6\cdot x}\ \Rightarrow\ \frac{1800}{144}=\frac{2400}{24x}\ \Rightarrow\ \frac{25}{2}=\frac{100}{x}\ \Rightarrow\ x=8\text{ gün}\]
(Bir makine bir saatte \(1800\div 144=12{,}5\) parça üretir; \(4\cdot 6\cdot 12{,}5=300\) parça/gün, \(2400\div 300=8\).)
\end{cozum}

\begin{alistirma}
\(25\) öğrencilik bir sınıfın not ortalaması \(64\) olarak hesaplanmıştır. Daha sonra bir öğrencinin \(82\) olan notunun yanlışlıkla \(32\) yazıldığı fark ediliyor. Sınıfın doğru ortalaması kaçtır?
\end{alistirma}
\begin{cozum}
Hatalı toplam: \(25\cdot 64=1600\). Doğru toplam: \(1600-32+82=1650\).\\
Doğru ortalama: \(\dfrac{1650}{25}=66\) \ (kısa yol: toplam \(50\) arttı, ortalama \(\frac{50}{25}=2\) artar).
\end{cozum}

\begin{alistirma}
İki pozitif sayının aritmetik ortalaması \(10\), geometrik ortalaması \(8\)'dir. Bu sayılar ve harmonik ortalamaları kaçtır?
\end{alistirma}
\begin{cozum}
\(a+b=20\), \(ab=64\): sayılar \(t^2-20t+64=(t-4)(t-16)=0\) kökleri, yani \(4\) ve \(16\).
\(\text{HO}=\frac{2ab}{a+b}=\frac{128}{20}=6{,}4\); \ \(10\ge 8\ge 6{,}4\) \checkmark
\end{cozum}

\kendinizoncesi
\section{Alıştırmalar (Kendiniz Çözün)}

\begin{enumerate}[label=\textbf{\arabic*.},leftmargin=*,itemsep=2pt]
\item \(\dfrac{x+5}{x-1}=\dfrac{7}{4}\) olduğuna göre \(x\) kaçtır?
\item \(\dfrac{a}{b}=\dfrac{3}{5}\) olduğuna göre \(\dfrac{a+b}{b-a}\) kaçtır?
\item \(\dfrac{a}{5}=\dfrac{b}{6}=\dfrac{c}{9}\) ve \(a+b+c=60\) olduğuna göre \(c\) kaçtır?
\item \(720\) TL, \(3\), \(4\) ve \(5\) ile doğru orantılı paylaştırılıyor. En az pay kaç TL'dir?
\item \(1:300\,000\) ölçekli bir haritada \(8\) cm olan uzaklık gerçekte kaç km'dir?
\item Birbirini çeviren iki dişli çarktan birinde \(40\), diğerinde \(60\) diş vardır. Küçük çark \(12\) tur attığında büyük çark kaç tur atar?
\item Eşit kapasiteli \(6\) musluk boş bir havuzu \(20\) saatte dolduruyor. Musluk sayısı \(2\) katına çıkarılırsa havuz kaç saatte dolar?
\item \(x\); \(y\) ile doğru, \(z\) ile ters orantılıdır. \(y=3\), \(z=4\) iken \(x=6\)'dır. \(y=5\), \(z=10\) iken \(x\) kaçtır?
\item \(5\) usta günde \(8\) saat çalışarak \(6\) günde \(120\) çift ayakkabı yapıyor. \(4\) usta günde \(10\) saat çalışarak \(160\) çift ayakkabıyı kaç günde yapar?
\item \(20\) kişilik bir grubun ortalaması \(55\), \(30\) kişilik bir grubun ortalaması \(65\)'tir. İki grubun birlikte ortalaması kaçtır?
\item \(5\) ve \(20\) sayılarının geometrik ve harmonik ortalamaları kaçtır?
\item Bir araç bir yolu saatte \(30\) km hızla gidip aynı yolu saatte \(70\) km hızla dönüyor. Ortalama hızı kaç km/sa'tir?
\item Ardışık \(11\) tek sayının toplamı \(231\)'dir. En büyüğü kaçtır?
\item \(a\) ile \(b\)'nin aritmetik ortalaması \(13\), geometrik ortalaması \(12\)'dir. \(a^2+b^2\) kaçtır?
\end{enumerate}

\needspace{8\baselineskip}
\subsection*{Cevap Anahtarı}
{\small
\begin{multicols}{2}
\begin{enumerate}[label=\textbf{\arabic*.},leftmargin=*,itemsep=1pt]
\item \(4x+20=7x-7\Rightarrow x=9\).
\item \(a=3k,\ b=5k\): \(\frac{8k}{2k}=4\).
\item \(20k=60\Rightarrow k=3\); \(c=27\).
\item \(12k=720\Rightarrow k=60\); \(3\cdot 60=180\) TL.
\item \(8\cdot 300\,000=2\,400\,000\) cm \(=24\) km.
\item \(40\cdot 12=60\cdot x\Rightarrow x=8\) tur.
\item \(6\cdot 20=12\cdot t\Rightarrow t=10\) saat.
\item \(k=\frac{6\cdot 4}{3}=8\); \(x=8\cdot\frac{5}{10}=4\).
\item \(\frac{120}{5\cdot 8\cdot 6}=\frac{160}{4\cdot 10\cdot x}\Rightarrow\frac12=\frac{4}{x}\Rightarrow x=8\) gün.
\item \(\frac{1100+1950}{50}=\frac{3050}{50}=61\).
\item \(\text{GO}=\sqrt{100}=10\); \(\text{HO}=\frac{200}{25}=8\).
\item \(\frac{2\cdot 30\cdot 70}{100}=42\) km/sa.
\item Ortanca \(\frac{231}{11}=21\); en büyük \(21+10=31\).
\item \(a+b=26,\ ab=144\): \(676-288=388\).
\end{enumerate}
\end{multicols}}


\vspace{2pt}
\noindent\textcolor{anamavi}{\rule{\textwidth}{0.5pt}}\\
{\small Gelecek hafta: Sayı, kesir ve yaş problemleri.}

\end{document}
